\uuid{B6pY}
\exo7id{5990}
\titre{exo7 5990}
\auteur{quinio}
\organisation{exo7}
\datecreate{2011-05-18}
\isIndication{false}
\isCorrection{true}
\chapitre{Probabilité discrète}
\sousChapitre{Probabilité et dénombrement}
\cpge{mpsi-mp2i.probabilites-univers-fini.e, pcsi.probabilites-variables-aleatoires.e, ptsi.probabilites-variables-aleatoires.e, mpsi-mp2i.probabilites-univers-fini.c, pcsi.probabilites-variables-aleatoires.c, ptsi.probabilites-variables-aleatoires.c, bcpst1.concepts-base-probabilites.a, bcpst1.concepts-base-probabilites.b, mp-mpi.variables-aleatoires-discretes.c, pc.probabilites-lois-usuelles.c, psi.probabilites-lois-usuelles.c, pt.probabilites-lois-usuelles.d}
\module{Probabilité et statistique}
\niveau{L2}
\difficulte{}

\contenu{
\question{
Dans un jeu de $52$ cartes, on prend une carte au hasard : les événements <<tirer un roi>> et 
<<tirer un pique>> sont-ils indépendants? quelle est la probabilité de <<tirer un roi ou
un pique>> ?
}
\reponse{
Soit $A$ : l'événement <<tirer un roi>>
et $B$ : <<tirer un pique>>.

$P(A\cap B)=\frac{1}{52};P(A)=\frac{4}{52}=\frac{1}{13};P(B)=\frac{13}{52}=\frac{1}{4}.$

Donc $P(A\cap B)=P(A)P(B)$ et donc les événements $A$ et $B$ sont indépendants.

$P(A\cup B)=P(A)+P(B)-P(A\cap B)=\frac{1}{13}+\frac{1}{4}-\frac{1}{52}=\frac{4}{13}$.
}
}
